% main.tex
% Copyright 2026 Zildjian E. California
%
% This file is part of the beamer-up derivative work and adapts the
% original beamer-uu sources by A.J.H. (Jos) Zuijderwijk to the
% maintained University of the Philippines `UP` theme surface.
%
% This work may be distributed and/or modified under the conditions of
% the LaTeX Project Public License, either version 1.3c of this license
% or (at your option) any later version. The latest version of this
% license is in
%   https://www.latex-project.org/lppl.txt
% and version 1.3c or later is part of all distributions of LaTeX
% version 2008 or later.
%
% This work has the LPPL maintenance status `maintained'.
%
% The Current Maintainer of this work is Zildjian E. California.
% Issue reports for this maintained derivative may be sent via
%   https://github.com/zcalifornia-ph/beamer-up/issues
%
% This work consists of the files beamercolorthemeUP.sty,
% beamerinnerthemeUP.sty, beamerouterthemeUP.sty, beamerthemeUP.sty,
% main.tex, and references.bib.
%
% Current checkout note: this deck instantiates the UP theme for the
% SSSP-modern CMSC 142 portfolio. This version establishes the deck
% spine and the four baseline-algorithm frames; later content passes
% fill the DMMSY centerpiece, experimental evaluation, author
% contributions, and timed dry-run refinements.
% Branding note: distributed UP-branded assets follow the UP branding
% guidelines and are intended for UP constituents/internal use. No
% official endorsement is implied.

\documentclass[aspectratio=169]{beamer}
\usetheme[showlogo,english]{UP}

% ===========================================
%  METADATA
% ===========================================
\title[SSSP: Dijkstra $\to$ DMMSY]{Single-Source Shortest Paths: From Dijkstra to DMMSY}
\subtitle{\textbf{A CMSC 142 portfolio study comparing five SSSP algorithms in pure-Python stdlib,\\with DMMSY (2025) as the centerpiece.}}
\author[California \& Rizal]{%
  Zildjian E. California\inst{1} \and
  Rey Marvin C. Rizal\inst{1}%
}
\institute{Department of Mathematics, Physics, and Computer Science\\College of Science and Mathematics\\University of the Philippines - Mindanao}
\addinstitute{1}{Department of Mathematics, Physics, and Computer Science,\\College of Science and Mathematics, University of the Philippines - Mindanao}
\venue{CMSC 142 Portfolio Presentation -- Assoc.\ Prof.\ Cinmayii G. Manliguez, Ph.D.}
\date{\today}
\multiauthor


\begin{document}
% ===========================================
%  TITLE SLIDE
% ===========================================
\begin{frame}[plain]
  \titlepage
\end{frame}
% Speaker note (California): name the paper centerpiece (DMMSY 2025) and
% the historical arc from Dijkstra 1959. Acknowledge Asc. Prof. Manliguez.
% Hand off to Rizal at "Five algorithms".

% ===========================================
%  OUTLINE
% ===========================================
\begin{frame}{Outline}
  \tableofcontents
\end{frame}
% Speaker note (California): point at the five sections; flag DMMSY as
% the centerpiece; flag Experimental Evaluation as the empirical claim.

% ===========================================
\section{Problem and Motivation}
% ===========================================

\begin{frame}{Single-Source Shortest Paths}
  \begin{block}{Problem}
    Given a weighted graph $G = (V, E)$ with edge weights $w \colon E \to \mathbb{R}$
    and a source vertex $s \in V$, compute
    $$d(v) = \min_{\pi : s \to v} \sum_{e \in \pi} w(e) \qquad \text{for every } v \in V.$$
  \end{block}

  \begin{itemize}
    \item Edge orientation: directed for Dijkstra, Bellman--Ford, A*, DMMSY; undirected for Thorup.
    \item Weight class: non-negative reals for Dijkstra / A* / DMMSY; signed reals for Bellman--Ford; non-negative integers for Thorup.
    \item Output: a distance map $d \colon V \to \mathbb{R}$ (and, optionally, a predecessor map for path reconstruction).
  \end{itemize}
\end{frame}
% Speaker note (California): one-sentence framing, then read the bullet
% list as the "ground rules" for the comparison. Do not derive the
% algorithm here -- that is the next section.

\begin{frame}{Five Algorithms Across the Historical Arc}
  \begin{itemize}
    \item \textbf{Dijkstra (1959)}~\cite{dijkstra1959note} -- the foundational non-negative-weight algorithm; binary-heap relaxation.
    \item \textbf{Bellman--Ford (1958)}~\cite{bellman1958routing}, drawing on Ford (1956)~\cite{ford1956network} -- handles negative weights; detects negative cycles.
    \item \textbf{A\textsuperscript{*} (1968)}~\cite{hart1968astar} -- heuristic-guided search; degenerates to Dijkstra at $h \equiv 0$.
    \item \textbf{Thorup (1999)}~\cite{thorup1999linear} -- linear-time undirected SSSP for non-negative integer weights in the word-RAM model.
    \item \textbf{DMMSY (2025)}~\cite{duan2025breaking} -- deterministic $O(m \log^{2/3} n)$ directed SSSP; \alert{the centerpiece of this study}.
  \end{itemize}

  \vskip 6pt
  We re-implement all five in pure-stdlib CPython 3.11+ and benchmark each against Dijkstra on a documented suite.
\end{frame}
% Speaker note (California): stress the arc and the centerpiece; flag
% that DMMSY is what breaks the sorting barrier. Hand off to Rizal.

% ===========================================
\section{Baseline Algorithms}
% ===========================================

\begin{frame}{Dijkstra (1959) -- Heap Relaxation}
  \begin{columns}[T]
    \begin{column}{0.52\textwidth}
      \begin{block}{Problem class}
        Non-negative weighted directed (or undirected) graphs; single source $s$.
      \end{block}
      \begin{block}{Design idea}
        Maintain a frontier set $B$ in a min-heap keyed by tentative distance; repeatedly extract the closest unsettled vertex, settle it, and relax its outgoing edges.
      \end{block}
    \end{column}
    \begin{column}{0.44\textwidth}
      \begin{exampleblock}{Complexity}
        Time $O((n+m) \log n)$; space $O(n)$.
      \end{exampleblock}
      \begin{alertblock}{Implementation}
        \texttt{src/sssp/dijkstra.py}; \texttt{heapq}-backed; 7 tests, 100\% line coverage.
        Source paper: \cite{dijkstra1959note}.
      \end{alertblock}
    \end{column}
  \end{columns}
\end{frame}
% Speaker note (California): "non-negative weights only" is the load-
% bearing assumption; this is why Bellman-Ford follows. Mention that
% Dijkstra is the BaseCase oracle for both DMMSY and the benchmark
% normalization.

\begin{frame}{Bellman--Ford (1958, Ford 1956) -- Edge Relaxation}
  \begin{columns}[T]
    \begin{column}{0.52\textwidth}
      \begin{block}{Problem class}
        Directed graphs with possibly negative edge weights; detect negative-weight cycles reachable from $s$.
      \end{block}
      \begin{block}{Design idea}
        Repeat full-graph edge relaxation $n-1$ times; one extra pass diagnoses any reachable negative cycle.
      \end{block}
    \end{column}
    \begin{column}{0.44\textwidth}
      \begin{exampleblock}{Complexity}
        Time $O(n m)$; space $O(n)$.
      \end{exampleblock}
      \begin{alertblock}{Implementation}
        \texttt{src/sssp/bellman\_ford.py}; reports \texttt{NegativeCycleReport}; 100\% line coverage.
        Source papers: \cite{bellman1958routing}, \cite{ford1956network}.
      \end{alertblock}
    \end{column}
  \end{columns}
\end{frame}
% Speaker note (Rizal): emphasize that BF is the "negative weights"
% baseline and is reused conceptually inside DMMSY's FindPivots. Note
% that the curated negative-cycle fixtures verify the detection branch.

\begin{frame}{A\textsuperscript{*} (1968) -- Heuristic Guidance}
  \begin{columns}[T]
    \begin{column}{0.52\textwidth}
      \begin{block}{Problem class}
        Goal-directed shortest path with a domain-specific heuristic $h(v)$ lower-bounding $d(v, \mathrm{goal})$.
      \end{block}
      \begin{block}{Design idea}
        Order the frontier by $f(v) = g(v) + h(v)$ instead of $g(v)$ alone; admissible $h$ preserves optimality.
      \end{block}
    \end{column}
    \begin{column}{0.44\textwidth}
      \begin{exampleblock}{Complexity}
        Time $O((n+m) \log n)$ worst case; expands fewer nodes than Dijkstra when $h$ is informative.
      \end{exampleblock}
      \begin{alertblock}{Implementation}
        \texttt{src/sssp/astar.py}, \texttt{src/sssp/heuristics.py}; ships \texttt{manhattan}, \texttt{zero}; \alert{$A^*(h\equiv 0) \equiv$ Dijkstra} tested.
        Source paper: \cite{hart1968astar}.
      \end{alertblock}
    \end{column}
  \end{columns}
\end{frame}
% Speaker note (California): the punchline is that A* generalizes
% Dijkstra; the zero-heuristic test makes that explicit. Mention the
% Manhattan grid demo as evidence of admissibility.

\begin{frame}{Thorup (1999) -- Linear-Time Undirected SSSP}
  \begin{columns}[T]
    \begin{column}{0.52\textwidth}
      \begin{block}{Problem class}
        Undirected graphs with non-negative integer edge weights; word-RAM model.
      \end{block}
      \begin{block}{Design idea}
        Hierarchical bucket structure exploiting integer weight bits and constant-time word operations to avoid the sorting bottleneck of comparison-based SSSP.
      \end{block}
    \end{column}
    \begin{column}{0.44\textwidth}
      \begin{exampleblock}{Complexity}
        Time $O(n + m)$ in the word-RAM model.
      \end{exampleblock}
      \begin{alertblock}{Implementation}
        \texttt{src/sssp/thorup99.py}; 11 tests, 99\% line coverage.
        \emph{Python $\ne$ word-RAM:} documented in ADR-002. Source paper: \cite{thorup1999linear}.
      \end{alertblock}
    \end{column}
  \end{columns}
\end{frame}
% Speaker note (Rizal): be explicit about the model gap -- Python's
% Big-Integer arithmetic does not preserve Thorup's linear-time
% guarantee. The implementation is faithful to the paper, not the
% asymptotic. This matters for the Experimental Evaluation discussion.

% ===========================================
\section{DMMSY (Centerpiece)}
% ===========================================

\begin{frame}{DMMSY (2025) -- Why It Matters}
  \small
  \begin{columns}[T]
    \begin{column}{0.52\textwidth}
      \begin{block}{Setting}
        Directed SSSP with non-negative real weights; labels use only comparison and addition.
      \end{block}
      \begin{block}{Barrier}
        Dijkstra-style methods still pay for global priority ordering, giving the sparse-graph comparison baseline $O(m + n \log n)$.
      \end{block}
    \end{column}
    \begin{column}{0.44\textwidth}
      \begin{exampleblock}{DMMSY claim}
        Duan--Mao--Mao--Shu--Yin break the barrier deterministically:
        $$O\!\left(m \log^{2/3} n\right).$$
      \end{exampleblock}
      \begin{alertblock}{Portfolio role}
        Centerpiece algorithm. Source paper: \cite{duan2025breaking}. Implementation: \texttt{src/sssp/dmmsy/}.
      \end{alertblock}
    \end{column}
  \end{columns}
\end{frame}
% Speaker note (California): make the model concrete: directed graph,
% non-negative real weights, only comparisons and additions. The reason
% this matters is that the paper breaks the sorting-style bottleneck
% that makes Dijkstra the historical baseline.

\begin{frame}{DMMSY Architecture}
  \footnotesize
  \begin{center}
  \resizebox{0.96\textwidth}{!}{%
  \begin{tikzpicture}[
    node distance=0.45cm and 0.45cm,
    box/.style={
      rectangle, rounded corners=2pt, minimum height=1.0cm,
      text width=2.8cm, align=center,
      font=\scriptsize\sffamily, draw=none
    },
    arr/.style={->, >=stealth, semithick, UPBlack80}
  ]
    \node[box, fill=UPMaroon, text=white] (transform) {Constant-degree\\transform\\{\tiny paper \S2}};
    \node[box, fill=UPForestGreen, text=white, right=of transform] (pivots) {FindPivots\\{\tiny Algorithm 1}};
    \node[box, fill=UPYellow, text=UPBlack, right=of pivots] (bmssp) {BMSSP\\recursion\\{\tiny Algorithms 2--3}};
    \node[box, fill=UPBlack, text=white, right=of bmssp] (blocklist) {BlockList\\frontier\\{\tiny Lemma 3.3}};

    \draw[arr] (transform) -- (pivots);
    \draw[arr] (pivots) -- (bmssp);
    \draw[arr] (bmssp) -- (blocklist);
    \draw[arr, bend left=25, UPForestGreen] (blocklist.north) to node[above, font=\tiny\sffamily, text=UPBlack] {bounded pulls} (bmssp.north);
  \end{tikzpicture}
  }
  \end{center}

  \begin{itemize}
    \item \textbf{Transform:} reduce high-degree vertices so work charges uniformly per arc.
    \item \textbf{FindPivots:} shrink the frontier to useful recursion roots.
    \item \textbf{BMSSP + BlockList:} solve bounded subproblems without globally sorting every label.
  \end{itemize}
\end{frame}
% Speaker note (California): walk left to right. The key message is not
% that one component is magic; the speedup comes from making frontier
% work local and bounded instead of globally sorted.

\begin{frame}{BMSSP Recursion}
  \scriptsize
  \begin{block}{Contract}
    \textsc{BMSSP}$(G,\ell,B,S,d,k,t)$ solves labels below boundary $B$ from source set $S$, then returns a smaller boundary $B'$ and settled set $U$.
  \end{block}

  \begin{columns}[T]
    \begin{column}{0.47\textwidth}
      \begin{exampleblock}{Level 0}
        \textsc{BaseCase}: bounded Dijkstra from one source, stopping after $k+1$ discoveries or boundary exhaustion.
      \end{exampleblock}
    \end{column}
    \begin{column}{0.49\textwidth}
      \begin{alertblock}{Level $>0$}
        \begin{enumerate}
          \item \textsc{FindPivots} shrinks $S$ to pivots $P$.
          \item BlockList pulls bounded batches below $B$.
          \item Recursive calls settle batches and relax outgoing edges.
        \end{enumerate}
      \end{alertblock}
    \end{column}
  \end{columns}

  \vskip 4pt
  \textbf{Key idea:} local bounded batches replace one global sorted frontier.
\end{frame}
% Speaker note (California): do not read this as code. Explain the
% recursion contract: each call solves only what is below a boundary,
% returns a separator, and lets the next call continue from there.

\begin{frame}{Parameters, Bound, and Implementation Caveat}
  \footnotesize
  \begin{columns}[T]
    \begin{column}{0.48\textwidth}
      \begin{exampleblock}{Paper parameterization}
        \[
          k = \left\lfloor (\log_2 n)^{1/3} \right\rfloor
        \]
        \[
          t = \left\lfloor (\log_2 n)^{2/3} \right\rfloor,\qquad
          \ell_{\max} = \left\lceil \frac{\log_2 n}{t} \right\rceil
        \]
        Balances pivot work, recursion depth, and block-list operations.
      \end{exampleblock}
    \end{column}
    \begin{column}{0.48\textwidth}
      \begin{block}{Asymptotic result}
        Deterministic directed SSSP:
        \[
          O\!\left(m \log^{2/3} n\right)\ \text{time}
        \]
        Space stays $O(n+m)$.
      \end{block}
      \begin{alertblock}{Python implementation boundary}
        CPython recursion overhead is real. Default large numeric benchmark runs use the documented fast path; \texttt{Weight} inputs and explicit parameters exercise recursive BMSSP.
      \end{alertblock}
    \end{column}
  \end{columns}
\end{frame}
% Speaker note (California): separate theory from implementation
% honestly. The paper result is the centerpiece; the portfolio proves a
% faithful implementation surface while acknowledging CPython overhead.

% ===========================================
\section{Experimental Evaluation}
% ===========================================

\begin{frame}{Experimental Evaluation -- To Be Filled}
  \begin{block}{Planned content}
    The experimental evaluation section will add:
    \begin{itemize}
      \item Methodology summary (warm runs, IO excluded, median + IQR; seed 2026; 60 rows).
      \item Per-algorithm $\times$ size-bucket median ratios vs.\ Dijkstra (table excerpt from \texttt{bench/results/sssp-benchmark-seed-2026.csv}).
      \item Embedded \texttt{bench/figures/runtime-by-algorithm.pdf} chart.
      \item NFR-F sanity: max DMMSY/Dijkstra ratio \texttt{1.093x} at $n \ge 1000$ (T-17 evidence).
    \end{itemize}
  \end{block}
\end{frame}
% Speaker note (Rizal): owns the empirical section; a later content pass will populate.

% ===========================================
\section{Reflection and Author Contributions}
% ===========================================

\begin{frame}{Reflection and Author Contributions -- To Be Filled}
  \begin{block}{Planned content}
    The closing content section will add:
    \begin{itemize}
      \item Honest framing of where the five algorithms are useful and where they break down (utility / limitations / DMMSY's Python caveat).
      \item AI Assistance Disclosure paragraph aligned with the U8 report (AC-08.3).
      \item Author Contributions slide naming California (1st) and Rizal (2nd) per US-09 / AC-09.2.
    \end{itemize}
  \end{block}
\end{frame}
% Speaker note (joint): California summarises usefulness/limitations;
% Rizal closes with Author Contributions and hands to Q&A.

% -- Standout --------------------------
\standoutframe{Honor. Excellence. Service.}

% -- Closing -----------------------------------
\thankframe{Questions?}{%
  Zildjian E. California \quad zecalifornia.com\\
  Rey Marvin C. Rizal\\
  Repository: \url{https://github.com/zcalifornia-ph/sssp-modern}%
}
% Speaker note (joint): leave the thankframe up during Q&A; both
% authors field questions; California fields DMMSY questions, Rizal
% fields experimental-evaluation questions.

% -- references -----------------------------
\begin{frame}[allowframebreaks]{References}
  The five source papers cited above are listed below; the Beamer class~\cite{tantau2004} and the UP Visual Identity Guidebook~\cite{upvig2017} are credited for the deck infrastructure.

  \vskip 6pt
  \printbibliography
\end{frame}

\end{document}
